\documentclass[12pt,compress,aspectratio=169]{beamer}
\input{../../mybeamer}

\usetikzlibrary{patterns}

\title{Classes 9 to 11: Rotational Motion of Rigid Bodies}
\subtitle{AP Physics C--Mechanics}
\input{../../me}
\input{../../mycommands}

\begin{document}

\begin{frame}
  \maketitle
\end{frame}


\begin{frame}{}
  This set of slides is the longest set in the entire course. It covers all the
  topics to be covered for the next 3 classes.
\end{frame}




\section{Introduction}

\begin{frame}{Uniform Circular Motion}
  Consider the uniform circular motion of an object with angular velocity
  $\vec\omega$.\footnote{If rotation is counterclockwise, $\vec\omega$ is
  \emph{out of the page}; if rotation is clockwise, $\vec\omega$ is
  \emph{into the page}. Or we can treat it as 1D motion: ($+$) if CCW; ($-$) if
  CW.}
  \begin{columns}
    \column{.3\textwidth}
    \centering
    \begin{tikzpicture}[scale=.7]
      \draw[axes] (-3,0)--(3,0) node[right]{$x$};
      \draw[axes] (0,-3)--(0,3) node[above]{$y$};
      \draw[dash dot] circle (2.5);
      \begin{scope}[rotate=38]
        \draw[vectors] (2.45,0)--(1,0) node[left]{$\vec F_c$};
        \draw[vectors] (2.5,0) arc (0:30:2.5) node[left]{$\omega$};
        \draw[mass] (2.5,0) circle (.1);
      \end{scope}
    \end{tikzpicture}

    \column{.7\textwidth}
    \begin{itemize}
    \item Centripetal force $\vec F_c\perp\vec v$
      \begin{itemize}
      \item $\vec F_c$ does not do any mechanical work
      \item Angular velocity $\vec\omega$ remains constant
      \end{itemize}
    \item\textbf{The ``rotational state'' of the object does not change}
    \item Rotation of an object is not determined by merely what forces are
      acting it
    \end{itemize}
  \end{columns}
  \vspace{.3in}
\end{frame}



\begin{frame}{Turning A Wrench}
  \begin{columns}
    \column{.3\textwidth}
    \pic{1.1}{wrench}
    
    \column{.7\textwidth}
    Similarly, when tightening/loosening a nut by turning a wrench,
    \begin{itemize}
    \item When the nut and the wrench turn, their ``rotational state'' change,
      i.e.\ they go from not rotating to rotating
    \item The applied force must be directed at some distance away from the
      nut
    \item How easy to turn the nut depends on \emph{both} the distance and the
      force
    \end{itemize}
  \end{columns}
\end{frame}



\begin{frame}{What is Torque?}
  \textbf{Torque} (also known as the \textbf{moment of force}, or just
  \textbf{moment}) is the tendency for a force to change the rotational state
  of a body.
  \begin{itemize}
  \item A force $\vec F$ acts at a point some position $\vec r$ from a
    \textbf{pivot}, or \textbf{fulcrum}
  %\item There is an angle $\phi$ between $\vec F$ and $\vec r$
  \item In the example below (may be a door handle), a force $\vec F$ is
    applied $\vec r$ away from the pivot at an angle $\phi$. A torque is
    generated around the pivot
  \end{itemize}
  \begin{center}
    \begin{tikzpicture}[scale=1.2]
      \draw[thick,fill=lightgray] (-.13,-.13) rectangle (5.13,.13);
      \draw[vectors,red] (0,0)--(5,0) node[midway,below]{$\vec r$};
      \draw[vectors,blue] (5,0)--(5.7,-1.5) node[right]{$\vec F$};
      \draw[dashed] (5,0)--(6,0);
      \draw[axes] (5.7,0) arc (0:-atan(1.5/.7):.7)
      node[midway,right]{$\phi$};
      \fill circle (.06) node[above=1]{\textbf{pivot}};
    \end{tikzpicture}
  \end{center}
\end{frame}



\begin{frame}{Torque}
  Torque $\vec\tau$ is defined as the cross product of the force $\vec F$ and
  the position vector $\vec r$ from the pivot, and where the force is applied.
  The SI unit for torque is a \textbf{newton metre} (\si{\newton\metre})
  
  \eq{-.1in}{
    \boxed{\vec\tau=\vec r\times\vec F}
  }
  \begin{itemize}
  \item\vspace{-.1in}The direction of the torque vector is orthogonal to both
    $\vec r$ and $\vec F$.
  \item The magnitude of torque:

    \eq{-.2in}{
      \boxed{\tau=rF_\perp=rF\sin\phi}
    }
    
    \vspace{-.07in}where $F_\perp$ is the component of $\vec F$ orthogonal
    (perpendicular) to the position vector $\vec r$
  \end{itemize}

  \vspace{.1in}{\footnotesize
    It is unlikely that you will do an actual cross product in an AP exam (it
    takes too long, and it does not actually test your knowledge in the
    physics), but you will need to know the basic properties of how a cross
    product applies in this situation.\par
  }
\end{frame}



\begin{frame}{Generating a Torque}
  \begin{columns}[T]
    \column{.33\textwidth}

    \vspace{-.09in}
    \begin{tikzpicture}[scale=.6]
      \draw[thick,fill=lightgray] (-.2,-.2) rectangle(5.2,.2);
      \draw[vectors,red] (0,0)--(5,0) node[midway,below]{$\vec r$};
      \draw[vectors,blue] (5,0)--(5,2) node[above]{$\vec F$};
      \fill circle (.08) node[above]{pivot};
    \end{tikzpicture}

    \vspace{.23in}
    {\footnotesize Force is perpendicular to the position vector, i.e.\
      $\vec r\perp\vec F$. In this case, the torque generated from the force is
      easy to calculate. The position vector $\vec r$ is known as the
      \textbf{moment arm} or \textbf{lever arm}.\par}
    
    \eq{-.2in}{
      \boxed{ \tau = rF }
    }

    \column{.33\textwidth}
    \begin{tikzpicture}[scale=.6]
      \draw[thick,fill=lightgray] (-.2,-.2) rectangle(5.2,.2);
      \draw[vectors,red] (0,0)--(5,0) node[midway,below]{$\vec r$};
      \draw[vectors,blue] (5,0)--(6.2,1.6) node[above]{$\vec F$};
      \draw[vectors,blue,dashed] (5,0)--(5,1.6) node[left]{$\vec F_\perp$};
      \draw[vectors,blue,dashed] (5,0)--(6.2,0) node[right]{$\vec F_\parallel$};
      \fill circle (.08) node[above]{pivot};
      \draw (5.6,0) arc (0:atan(4/3):.6) node[pos=2/3,right=0]{$\phi$};
    \end{tikzpicture}

    \vspace{.23in}{\footnotesize Force is applied at an angle $\phi$ to the
      position vector. Only the perpendicular component ($F_\perp$) generates a
      torque; the parallel component ($F_\parallel$) does not.\par}
    
    \eq{-.23in}{
      \boxed{
        \tau = r{\color{blue}F_\perp}=r{\color{blue}F\sin\phi}
      }
    }
    
    \column{.33\textwidth}
    \vspace{-.01in}
    \begin{tikzpicture}[scale=.6]
      \draw[thick,fill=lightgray] (-.2,-.2) rectangle(5.2,.2);
      \draw[vectors,red] (0,0)--(5,0) node[midway,below]{$\vec r$};
      \draw[vectors,blue] (5,0)--(6.2,1.6) node[above]{$\vec F$};
      \begin{scope}[rotate=-atan(3/4)]
        \draw[vectors,red] (0,0)--(4,0) node[midway,below]{$\vec r_\perp$};
        \draw (4,-.5)--(4,5.5)
        node[pos=0,right]{\small line of action of $\vec F$};
        \draw (0,-1)--(0,2);
        \draw[vectors,blue,dash dot] (4,0)--(4,2) node[midway,right]{$\vec F$};
        \fill[blue] (4,0) circle (.04);
      \end{scope}
      \fill circle (.08) node[above]{pivot};
      \draw (5,0)--(6.3,0);
      \draw (5.6,0) arc (0:atan(4/3):.6) node[pos=2/3,right]{$\phi$};
      \draw (.6,0) arc (0:atan(4/3):.6) node[pos=2/3,right]{$\phi$};
    \end{tikzpicture}

    {\footnotesize $\vec F$ can be moved along its line of action until it
      intersects the perpendicular component of $\vec r$\par}
    
    \eq{-.24in}{
      \boxed{
        \tau = {\color{red}r_\perp}F
        ={\color{red}r}F{\color{red}\sin\phi}
      }
    }

    \vspace{-.25in}{\footnotesize The moment arm is $\vec r_\perp$ in this
      case.}
  \end{columns}
\end{frame}



\begin{frame}{\emph{NOT} Generating a Torque}
  No torque is generated if the line of action of the force goes through the
  pivot.
  \begin{columns}[T]
    \column{.5\textwidth}
    \centering
    \vspace{.52in}
    \begin{tikzpicture}[scale=.7]
      \draw[thick,fill=lightgray] (-.2,-.2) rectangle(5.2,.2);
      \draw[vectors,red] (0,0)--(5,0) node[midway,below]{$\vec r$};
      \draw[vectors,blue] (5,0)--(7,0) node[above]{$\vec F$};
      \draw (-1,0)--(9,0) node[above]{line of action};
      \fill circle (.08) node[above]{pivot};
      \fill[blue] (5,0) circle (.04);
    \end{tikzpicture}

    \vspace{.03in}{\footnotesize
      Force $\vec F$ is parallel, or anti-parallel, to $\vec r$ (i.e.\
      $F_\perp=0$).\par}

    \column{.5\textwidth}
    \centering
    \begin{tikzpicture}[scale=.7]
      \draw[thick,fill=lightgray] (-.2,-.2) rectangle(5.2,.2);
      \begin{scope}[rotate=-atan(3/4)]
        \draw (0,-1)--(0,3) node[right]{line of action};
        \draw[vectors,blue] (0,0)--(0,2) node[right]{$\vec F$};
      \end{scope}
      \fill circle (.08) node[above]{pivot};
    \end{tikzpicture}

    {\footnotesize
      $\vec F$ is applied at the pivot (i.e.\ $\vec r=\vec 0$).\par
    }
  \end{columns}
\end{frame}



\begin{frame}{Net Torque}
  Net torque is the vector sum of all torques acting at a point. Keep in mind
  the sign convention:

  \eq{-.1in}{
    \boxed{
      \vec\tau_\text{net}=\sum_{i=1}^N\vec\tau_i
    }
  }
  \begin{itemize}
  \item Torque can come from individual forces acting at a single point:

    \eq{-.15in}{ \tau_i = r_iF_i\sin\phi_i }

  \item\vspace{-.15in}Often the forces that generate torques are spread
    over an area, e.g.\ force on the head of a screwdriver
  \item For rotational motion, we only care about the torque; generally
    we don't have to worry about exactly how the force generates the torque
  \end{itemize}
\end{frame}


\begin{frame}{Net Torque}
  
\end{frame}


\begin{frame}{Example Problem}
  \textbf{Example:} Find the net torque on point C.
  \begin{center}
    \begin{tikzpicture}[scale=2.65]
      \begin{scope}[thick]
        \fill[blue!40,draw=black] (-1.65,-.12) rectangle (1.65,.12);
        \draw[<->] (-1.5,-.25)--(1.5,-.25) node[midway,fill=black!2]{3.0 m};
        \draw[<->] (-1.5,.25)--(0,.25) node[midway,fill=black!2]{1.5 m};
        \draw[dashed] (-2.2,0)--(2.2,0);
        \draw[dashed] (0,0)--(0,.5);
      \end{scope}
      \begin{scope}[orange,vectors]
        \draw[rotate=-45] (0,0)--(0,1.5) node[right]{\SI{30}\newton};
        \draw[rotate around={30:(-1.5,0)}] (-1.5,0)--(-2.5,0)
        node[left]{\SI{20}\newton};
        \draw[rotate around={-30:(1.5,0)}] (1.5,0)--(2.2,0)
        node[right]{\SI{10}\newton};
      \end{scope}
      \draw[axes] (0,.4) arc (90:45:.4) node[pos=.7,above]{\ang{45}};
      \draw[axes] (-1.9,0) arc (180:210:.4) node[midway,left]{\ang{30}};
      \draw[axes] (1.9,0) arc (0:-30:.4) node[midway,right]{\ang{30}};
      \fill (-1.5,0) circle (.1) node[white]{\textbf{A}};
      \fill circle (.1) node[white]{\textbf{B}};
      \fill (1.5,0) circle (.1) node[white]{\textbf{C}};
    \end{tikzpicture}
  \end{center}
\end{frame}



\begin{frame}{Example Problem}
  \begin{columns}
    \column{.75\textwidth}
    Revisiting a problem that we encountered in Class \#2 (dynamics):

    \vspace{.2in}\textbf{Example:} A uniform ladder is \SI{5.0}{\metre} long
    and weighs \SI{400}\newton. The ladder rests against a slippery vertical
    wall, as shown in the figure. The inclination angle between the ladder and
    the rough floor is \ang{53}. Find the reaction forces from the floor and
    from the wall on the ladder and the coefficient of static friction $\mu_s$
    at the interface of the ladder with the floor that prevents the ladder from
    slipping.

    \column{.25\textwidth}
    \pic1{../Class02b-dynamics/graphics/ladder}
  \end{columns}
\end{frame}



\section{Angular Momentum}

\begin{frame}{Angular Momentum}
  Consider a mass $m$ connected to a massless beam rotates with velocity
  $\vec v$ at a position $\vec r$ from the centre (shown on the right). It has
  an \textbf{angular momentum} ($\vec L$), defined as:
  \begin{columns}
    \column{.77\textwidth}
    
    \eq{-.1in}{
      \boxed{\vec L=\vec r\times\vec p=m(\vec r\times\vec v)}
    }

    \vspace{-.1in}Expanding the term with $\vec v=\vec\omega\times\vec r$, we
    can express angular momentum only in quantities related to rotations:

    \eq{-.1in}{
      \boxed{\vec L=m(\vec r\times\vec v)
        =m(\vec r\times(\vec\omega\times\vec r))
        =mr^2\vec\omega}
    }
    
    \vspace{-.2in}Or in scalar form:
    
    \eq{-.1in}{
      \boxed{L=rmv=mr^2\omega}
    }

    \vspace{-.2in}The unit for angular momentum is a
    \textbf{kilogram metre squared per second}
    (\si{\newton\metre\squared\per\second}).
    
    \column{.23\textwidth}
    \begin{tikzpicture}[scale=2.5]
      \begin{scope}[rotate=70]
        \draw[line width=3.5] (0,0)--(2,0);
        \fill (2,0) circle (.04) node[right]{$m$};
        \fill circle (.05) node[right]{pivot};
        \draw[vectors,red] (0,0)--(2,0) node[midway,right]{$\vec r$};
        \draw[vectors,cyan] (2,0)--(2,.5) node[left]{$\vec v$};
      \end{scope}
    \end{tikzpicture}
  \end{columns}
\end{frame}



\section{Moment of Inertia}

\begin{frame}{Moment of Inertia}
  %Look again at the definition of angular momentum:
    
  \eq{-.1in}{
    \vec L=\underbracket[1pt]{mr^2}_I\vec\omega
  }
    
  The quantity $I=mr^2$ is called the \textbf{moment of inertia}\footnote{Also
  known as \textbf{rotational inertia} in AP Physics in recently years}, with
  an SI unit of \textbf{kilogram metre squared} (\si{\kilo\gram\metre\squared}).

  \eq{-.1in}{
    \boxed{\vec L=I\vec\omega}
  }

  \vspace{-.15in}Moment of inertia can be considered to be an object's
  ``rotational mass'' (i.e.\ moment of inertia in rotational motion is
  equivalent to what mass is for translational motion)
  \vspace{.3in}
\end{frame}



\begin{frame}{Moment of Inertia}
  For a \emph{single particle} of mass $m$ rotating at a distance $r$ from the
  pivot:
  
  \eq{-.1in}{
    \boxed{I=mr^2}
  }

  \vspace{-.1in}For $N$ \emph{discrete} particles rotating at the same
  $\omega$, each mass $m_i$ at distance $r_i$ from the pivot:

  \eq{-.2in}{
    \boxed{I=\sum_{i=1}^N m_ir_i^2}
  }

  For a \emph{continuous distribution of mass} (i.e.\ as $N\to\infty$) rotating
  about a pivot, integral calculus is need:

  \eq{-.2in}{
    \boxed{I=\int r^2\dl m}
  }
\end{frame}



\begin{frame}{Moment of Inertia}
  \centering
  \pic{.6}{mic}\\
  In AP Physics C, you may be asked to derive the moment of inertia for simple
  shapes.
\end{frame}


\begin{frame}[t]{Example Problem}
  \textbf{Example:} Calculate the rotational inertia about the end of a thin
  rod of length $L$ with constant linear mass density $\gamma$:
  \begin{center}
    \begin{tikzpicture}
      \draw[mass] rectangle (10,.2);
      \draw[|<->|,thick] (0,.4)--+(10,0) node[midway,above]{$L$};
      \fill (0,.1) circle (.07) node[midway,below]{pivot};
    \end{tikzpicture}
  \end{center}
\end{frame}



\begin{frame}[t]{Example Problem}
  \textbf{Example:} Calculate the rotational inertia about the centre of mass
  of a thin rod of length $L$ with constant linear mass density $\gamma$:
  \begin{center}
    \begin{tikzpicture}
      \draw[mass] rectangle (10,.2);
      \draw[|<->|,thick] (0,.4)--+(10,0) node[midway,above]{$L$};
      \fill (5,.1) circle (.07) node[below=3]{pivot};
    \end{tikzpicture}
  \end{center}
\end{frame}



\begin{frame}{Parallel Axis Theorem}
  \begin{columns}
    \column{.3\textwidth}
    \centering
    \pic{.7}{parallel-axis}
    
    \column{.7\textwidth}
    The \textbf{parallel axis theorem} relates the moment of inertia of an
    object along two different but parallel axes by:

    \eq{-.1in}{
      \boxed{
        I=I_\text{cm}+md^2
      }
    }

    We can use the theorem to verify the results from the previous two examples.
  \end{columns}
\end{frame}




\begin{frame}{Angular Momentum and Moment of Inertia}
  Linear and angular momentum have very similar expressions
    
  \vspace{-.3in}{\large
    \begin{align*}
      \vec p &= m\vec v\\
      \vec L &= I\vec\omega
    \end{align*}
  }
  
  Just as $\vec p$ describes the \emph{translational} state of motion of an
  object, $\vec L$ describes its \emph{rotational} state
\end{frame}



\section{Laws of Motion}

\begin{frame}{Second Law of Motion for Rotational Motion}
  The net torque is the time rate of change of angular momentum:

  \eq{-.1in}{
    \vec\tau_\text{net}=\vec r\times\vec F_\text{net}
    =\vec r\times\diff{\vec p}t
    =\diff{(\vec r\times\vec p)}t\;\;\longrightarrow\;\;
    \boxed{\vec\tau_\text{net}=\diff{\vec L}t}
  }
  \begin{itemize}
  \item If the net torque on a system is zero, then the rate of change
    of angular momentum is zero, and we say that the angular momentum is
    conserved. 
  \item e.g.\ When an ice skater starts to spin and draws his arms inward.
    Since angular momentum is conserved, a decrease in $r$ means an increase in
    $\omega$. The mass of the skater cannot change, but the distance of the mass
    from the rotational axis can
  \end{itemize}
\end{frame}


\begin{frame}{Equilibrium: First Law of Motion}
  \begin{columns}
    \column{.5\textwidth}
    An object is in \textbf{translational equilibrium}
    when the net \emph{force} acting on it is zero:
  
    \eq{-.1in}{
      \vec F_\text{net}=\vec 0
    }

    % This does \emph{not} mean that the object has no translational motion;
    %it just means that
    The object's translational state is not changing, i.e.\
    momentum $\vec p$ is constant. For constant mass $m$, this means that
    $\vec a_\text{cm}=\vec 0$ and $\vec v_\text{cm}=\text{constant}$
    
    \column{.5\textwidth}
    An object is in \textbf{rotational equilibrium} when
    the net \emph{torque} acting on it is zero:

    \eq{-.1in}{
      \vec\tau_\text{net}=\vec 0
    }

    %This does \emph{not} mean that the object has no rotational motion; it just
    %means that
    The object's rotational state is not changing, i.e.\
    angular momentum $\vec L$ is constant. For constant moment of inertia $I$,
    this means that $\vec\alpha=\vec 0$, and $\vec\omega=\text{constant}$
  \end{columns}
\end{frame}



%\begin{frame}{Equilibrium: First Law of Motion}
%  An object is in \textbf{translational equilibrium} is when the net force
%  acting on it is zero:
%  
%  \eq{-.1in}{
%    \vec F_\text{net}=\vec 0
%  }
%
%  This does \emph{not} mean that the object has no translational motion; it
%  just means that the object's translational state is not changing, i.e.\
%  momentum $\vec p$ is constant. For constant mass $m$, this means that
%  $\vec a=\vec 0$ and $\vec v=\text{constant}$.
%\end{frame}
%
%
%
%
%\begin{frame}{Equilibrium: First Law of Motion}
%  Likewise, an object is in \textbf{rotational equilibrium} when the net torque
%  acting on it is zero:
%
%  \eq{-.1in}{
%    \vec\tau_\text{net}=\vec 0
%  }
%  
%  This does \emph{not} mean that the object has no rotational motion; it just
%  means that the object's rotational state is not changing, i.e.\
%  angular momentum $\vec L$ is constant. For constant moment of inertia $I$,
%  this means that $\vec\alpha=\vec 0$, and $\vec\omega=\text{constant}$.
%\end{frame}



\begin{frame}{Second Law of Motion}% for Translational Motion}
  \begin{columns}[T]
    \column{.5\textwidth}
    For translational motion, the %general form of the first and
    second law of motion relate the net force to the rate of change of the
    object's momentum:

    \eq{.02in}{
      \boxed{
        \vec F_\text{net}=\diff{\vec p}t
      }
    }

    For objects with constant mass, this reduces to the more familiar form:

    \eq{-.1in}{
      \boxed{
        \vec F_\text{net}=m\vec a_\text{cm}
      }
    }

    \column{.5\textwidth}
    %\end{frame}
    %
    %
    %
    %\begin{frame}{First \& Second Law of Motion for Rotational Motion}
    Likewise, for rotational motion, the second law of motion relate the net
    torque to the rate of change of the object's angular momentum:
    %Likewise, the equation for the first and second laws of motion for
    %rotational motion has a similar form, but with net torque
    %$\vec\tau_\text{net}$ replacing net force $\vec F_\text{net}$, and angular
    %momentum $\vec L$ replacing linear momentum $\vec p$:

    \eq{-.2in}{
      \boxed{
        \vec\tau_\text{net}=\diff{\vec L}t
      }
    }

    For objects with constant momentum of inertia $I$, this reduces to:
    
    \eq{-.1in}{
      \boxed{
        \vec\tau_\text{net}=I\vec\alpha
      }
    }
  \end{columns}
\end{frame}



\begin{frame}{But there is no rotational motion, is there?}
  Even when there is no apparent rotational motion, it does not necessarily
  mean that angular momentum is zero! In this case, mass $m$ travels along a
  straight path at constant velocity (uniform motion), but the angular momentum
  around point $P$ is \emph{not} zero:
  \begin{center}
    \begin{tikzpicture}
      \draw[dashed] (-5,0)--(5,0);
      \draw[vectors,red] (-5,0)--(-3,0) node[below]{$\vec v$};
      \shade[ball color=red] (-5,0) circle (.15) node[above=4]{$m$};
      \draw[vectors,blue] (0,-2)--(-5,0)
      node[midway,below,rotate=-atan{2/5}]{$\vec r_1$};

      \draw circle (.15) node[above=4]{$m$};
      \draw[vectors,red!40] (.2,0)--(2,0) node[right]{$\vec v$};
      \draw[vectors,blue!40] (0,-2)--(0,0) node[midway,left]{$\vec r_2$};

      \draw (4,0) circle (.15) node[above=4]{$m$};
      \draw[vectors,red!40] (4.2,0)--(6,0) node[right]{$\vec v$};
      \draw[vectors,blue!40] (0,-2)--(4,0)
      node[midway,below,rotate=atan{2/4}]{$\vec r_3$};

      \fill (0,-2) circle (.05) node[below]{$P$};
    \end{tikzpicture}
  \end{center}
  There is no force and no torque acting on the object, therefore both linear
  momentum ($\vec p=m\vec v$) and angular momentum
  ($\vec L=\vec r\times\vec p$) are constant
\end{frame}



\begin{frame}{Example Problem}
  \textbf{Example:} A skater extends her arms (both arms!), holding a
  \SI{2.0}{\kilo\gram} mass in each hand. She is rotating about a vertical axis
  at a given rate. She brings her arms inward toward her body in such a way that
  the distance of each mass from the axis changes from \SI{1.0}{\metre} to
  \SI{.50}\metre. Her rate of rotation (neglecting her own mass) will?
\end{frame}




\begin{frame}{Example Problem}
  \textbf{Example:} A \SI{1.0}{\kilo\gram} mass swings in a vertical circle
  after having been released from a horizontal position with zero initial
  velocity. The mass is attached to a massless rigid rod of length
  \SI{1.5}\metre. What is the angular momentum of the mass, when it is in its
  lowest position?
\end{frame}



\section{Work \& Energy in Rotational Motion}

\begin{frame}{Mechanical Work}
  For translational motion\footnote{Showing 1D motion for simplicity},
  mechanical work is defined as

  \eq{-.1in}{
    W_\text{trans}=\int_{x_0}^{x_1}F\dl x
  }

  For rotational motion, mechanical work is defined similarly as:

  \eq{-.1in}{
    W_\text{rot}=\int_{x_0}^{x_1}F\dl x=\int_{\theta_0}^{\theta_1}F(r\dl\theta)
    \quad\to\quad
    \boxed{
      W_\text{rot}=\int_{\theta_0}^{\theta_1}\tau\dl\theta
    }
  }
  \vspace{.2in}

  Like in deriving the laws of motion, all we have done is to replace
  $F\to\tau$, and $\dl x\to\dl\theta$
\end{frame}



\begin{frame}{Rotational Kinetic Energy}
  When a net \emph{torque} causes an angular acceleration in a rotational
  system (assuming that the system has a constant \emph{moment of inertia}),
  the total (net) rotational work done on the object $W_\text{net}$ is given by:

  \eq{-.1in}{
    W_\text{net,rot}
    =\int_{\theta_0}^{\theta_1} \tau_\text{net}(\theta)\dl\theta
    =\int I\alpha\dl\theta
    =I\int\diff{\omega}t\dl\theta
  }

  Since both $\omega(t)$ and $\theta(t)$ are continuously differentiable in
  time, we can switch the order of differentiation:
  
  \eq{-.1in}{
    =I\int\diff{\theta}t\dl\omega=I\int_{\omega_0}^{\omega_1}\omega\dl\omega
  }
  where $\omega_0=\omega(\theta_0)$ and $\omega_1=\omega(\theta_1)$.
\end{frame}



\begin{frame}{Rotational Kinetic Energy}
  This integral, when integrated from $\omega_0$ to $\omega_1$, becomes:

  \eq{-.1in}{
    =I\int_{\omega_0}^{\omega_1}\omega\dl\omega
    =\frac12I\omega^2\Big|^{\omega_1}_{\omega_0}
    =\frac12I\omega_1^2-\frac12I\omega_0^2
    =\Delta K_\text{rot}
  }
  
  where $K_\text{rot}$ is defined as the \textbf{rotational kinetic energy}:

  \eq{-.1in}{
    \boxed{
      K_\text{rot}=\frac12I\omega^2
    }
  }
  This is, of course, the same mathematics that we did for translational work.
\end{frame}



\begin{frame}{Work-Energy Theorem}
  Like translational kinetic energy, the definition of \emph{rotational}
  kinetic energy came from integrating net work, in that net \emph{rotational}
  work equals to the change in \emph{rotational} kinetic energy:

  \eq{-.2in}{
    \boxed{
      W_\text{net,rot}=\Delta K_\text{rot}
    }
  }
  \begin{itemize}
  \item If net rotational work is positive ($W_\text{net,rot}>0$)
    \begin{itemize}
    \item $K_\text{rot}$ increases ($\Delta K_\text{rot}>0$) by the same amount
    \item The object's angular speed increases
    \end{itemize}
  \item If net rotational work is negative ($W_\text{net,rot}<0$)
    \begin{itemize}
    \item $K_\text{rot}$ decreases ($\Delta K_\text{rot}<0$) by the same amount
    \item The object's angular speed decreases
    \end{itemize}
  \item It does not matter \emph{what} net torque consists of
  \end{itemize}
\end{frame}




\begin{frame}{Rotational Kinetic Energy}
  Alternatively, to find the kinetic energy of a rotating system of particles
  (discrete number of particles, or continuous mass distribution), we also
  sum/integrate the kinetic energies of the individual particles:
    
  \eq{-.1in}{
    K_\text{rot}=\sum_{i=1}^N\frac12m_iv_i^2=\sum_{i=1}^N\frac12m_i(r_i\omega)^2
    =\frac12\left[\sum_{i=1}^N m_ir_i^2\right]\omega^2 = \frac12I\omega^2
  }
\end{frame}



\begin{frame}{Kinetic Energy of a Rotating System}
  The total kinetic energy of a rotating system is the sum of its translational
  and rotational kinetic energies at its centre of mass:

  \eq{-.1in}{
    \boxed{
      K_\text{tot}
      =K_\text{trans}+K_\text{rot}
      =\frac12mv_\text{cm}^2+\frac12I_\text{cm}\omega^2
    }
  }
  
  In this case, $I_\text{cm}$ is calculated at the CM. For simple problems, we
  only need to compute rotational kinetic energy at the pivot:

  \eq{-.1in}{
    \boxed{K_\text{tot}=\frac12I_\text{p}\omega^2}
  }
  
  In this case, the $I_\text{p}$ is calculated at the pivot.
  \textbf{IMPORTANT:} $I_\text{cm}\neq I_\text{p}$
\end{frame}



\begin{frame}{Power}
  As you have seen in other physics course (and earlier in this course),
  \textbf{power} is the \emph{rate} at which work is done, i.e.\ the rate at
  which energy is being transformed. A rotational system is defined the same
  way as it was in the translational system:
  
  \eq{-.1in}{
    \boxed{P(t) = \diff Wt}\quad\quad
    \boxed{\overline P = \frac W{\Delta t}}
  }
\end{frame}



\begin{frame}{Power}
  If a torque is used to push an object, the instantaneous power produced by the
  force is:
  
  \eq{-.1in}{
    P(t)=\diff Wt=\frac{\tau\dl\theta}{\dl t}
    =\tau\diff{\theta}t \quad\to\quad
    \boxed{P(t)=\tau(t)\omega(t)}
  }
  
  Consider an engine that is able to produce a torque $\tau$, the total power
  that it generates (i.e.\ converts from fuel) is equal to $\tau$ multiplied
  by how fast it is spinning (i.e. $\omega$):
\end{frame}




\section{Momentum in a Rotational System}

\begin{frame}{Conservation of Momentum}
  During a collision/explosion, there is no external force, therefore the total
  linear momentum of the system is conserved:

  \eq{-.1in}{
    \boxed{
      \sum\vec p_i=\sum\vec p_i'
    }
  }
  
  \vspace{-.1in}But if there is no external force, there is also be no external
  torque. Therefore, total \emph{angular} momentum of the system must also be
  conserved in a collision:
  
  \eq{-.1in}{
    \boxed{
      \sum\vec L_i=\sum\vec L_i'
    }
  }

  Additionally, if the collision is elastic, then kinetic energy (both
  translational and rotational) is also conserved.

  \eq{-.25in}{
    \boxed{
      \sum K_i=\sum K_i'
    }
  }
  
  %Let's remind ourselves the sign convention for angular momentum on the
  %$xy$-plane: $L$ is $+$ when rotational motion is counter clockwise, and $-$
  %when it is clockwise
\end{frame}





\begin{frame}{Spinning Disk - Pivot at the Centre of Mass}
%  \begin{columns}
%    \column{.3\textwidth}
  \begin{tikzpicture}
    \draw[mass] circle (2);
    \draw[thick] circle (.13);
    \fill (0,0)--(.13,0) arc (0:90:.13) node[above=-2]{cm/pivot}
    --(0,0);
    \fill (0,0)--(-.13,0) arc (180:270:.13)--(0,0);
    \draw[axes,rotate=-30] (1,0) arc (0:60:1) node[above]{$\omega$, $\alpha$};
    \draw[axes,rotate=40] (0,0)--(-2,0) node[midway,below=3]{$R$};
    \node at (0,-2.5) {Mass $M$};
    \node at (0,-3.3) {$I=\dfrac12MR^2$};
%    \only<2>{
%      \begin{scope}[violet]
%        \draw[thick,dash dot] (0,-.13)--(0,-2);
%        \foreach \y in {-.4,-.8,...,-2.1} \draw[vectors] (0,\y)--(-.8*\y,\y);
%        \node[right] at (1.6,-2) {$\vec v=\omega R$};
%      \end{scope}
%    }
  \end{tikzpicture}

%    \column{.7\textwidth}
%    \only<2>{Speed at any point on this disk is proportional to its distance
%      to the pivot:
%
%      \eq{-.1in}{
%        v=\omega r \quad\text{for}\quad 0\leq r\leq R
%      }
%    }
%    \only<3>{When it is spinning, it has an angular momentum of
%
%      \eq{-.1in}{
%        L = I\omega
%      }
%
%      and a rotational kinetic energy:
%      
%      \eq{-.1in}{
%        K_r = \frac12I\omega^2
%      }
%    }
%    
%    \item Any angular acceleration that it may have will be due to the net
%      torque at the pivot/CM:
%      \begin{displaymath}
%        \tau_\text{net}=\frac{\Delta L}{\Delta t}=I\alpha
%      \end{displaymath}
%    \end{itemize}
%  \end{columns}
\end{frame}



\begin{frame}{A Rotating Rod}
  \begin{columns}
    \column{.3\textwidth}
    \centering
    \begin{tikzpicture}
      \draw[mass] (-.15,0) rectangle (.15,-5);
      \draw[thick] (0,-2.5) circle (.13);
      \fill (0,-2.5)--(.13,-2.5) node[right=-1]{cm} arc(0:90:.13)
      --(0,-2.5);
      \fill (0,-2.5)--(-.13,-2.5) arc(180:270:.13)--(0,-2.5);
      \fill circle (.07) node[above]{pivot};
      \draw[axes,rotate=-90] (.7,0) arc (0:180:.7)
      node[left]{$\omega$, $\alpha$};
      \draw[thick,|<->|] (-.5,0)--(-.5,-5) node[midway,fill=black!2]{$L$};
%      \node at (0,-2.5) {Mass $M$};
%      \node at (0,-3.3) {$I=\dfrac12MR^2$};
%      \only<2>{
%        \begin{scope}[violet]
%          \draw[thick,dash dot] (0,-.13)--(0,-2);
%          \foreach \y in {-.4,-.8,...,-2.1}\draw[vectors] (0,\y)--(-.8*\y,\y);
%          \node[right] at (1.6,-2) {$\vec v=\omega R$};
%        \end{scope}
%      }
    \end{tikzpicture}

    \column{.7\textwidth}
  \end{columns}
\end{frame}


\begin{frame}{Solving Rotational Problems}
  When solving for rotational problems like the ones described in the previous
  sections:
  \begin{itemize}
  \item Draw a free-body diagram to account for all forces
  \item The direction of friction force is not always obvious
  \item The magnitude of any static friction force cannot be assumed to be at
    maximum.
  \item If the object is to change its rotational state, there must be a net
    torque causing it.
  \end{itemize}
\end{frame}



\begin{frame}{Solving Rotational Problems}
  Once the free-body diagram is complete, the forces should break down into
  their \emph{forces} into $\iii$, $\jjj$ and $\kkk$ components. If the axes
  are defined properly, only one direction should have acceleration (usually
  $\iii$), i.e.:
  
  \eq{-.1in}{
    \sum F_x=ma_\text{cm} \quad\quad \sum F_y=0 %\quad\quad \sum F_z=0
  }

  For AP Physics C, there will also be one equation for rotation, as torque is
  only applied in one direction:
    
  \eq{-.1in}{
    %\sum\tau_x=0 \quad\quad \sum\tau_y=0 \quad\quad
    \sum\tau=I\alpha
  }
\end{frame}



\begin{frame}{Solving Rotational Problems}
  For rotational motion dynamics equation:
  \begin{enumerate}
  \item Relate the force(s) that causes rotational motion to the net torque

    \eq{-.1in}{
      \tau_\text{net}=\sum_i r_iF_i
    }
  \item Substitute the expression for momentum of inertia (which has both mass
    and radius terms in it) into the equation for rotational motion
  \item Relate angular acceleration to linear acceleration, if applicable:

    \eq{-.1in}{
      \alpha=\frac{a_\text{cm}}R
    }
  \end{enumerate}
  Now there are two equations with force and acceleration terms.
\end{frame}




\section{Rolling Problems}


\begin{frame}{Motion with Both Translation and Rotation}
  Often motion of an object includes both translational and rotational motion.
  \begin{center}
    \begin{tikzpicture}
      \fill circle (1);
      \fill[black!2] circle (.8);
      \fill circle (.1);
      \foreach \theta in {0,30,...,359} \draw[thick,rotate=\theta] (0,0)--(1,0);
      \draw[thick] (-3,-1)--(3,-1);
      \draw[vectors,red] (1.1,0)--+(1,0) node[right]{$\vec v_\text{cm}$};
      \draw[vectors,red,rotate=40] (.3,0) arc (0:-140:.3) node[left]{$\omega$};
    \end{tikzpicture}
  \end{center}
  For example, a bicycle tire that is rolling on the road without slipping has
  both translational velocity $\vec v_\text{cm}$ at its centre of mass, and
  angular velocity $\omega$.
\end{frame}



\begin{frame}{No-Slip Rolling}
  \begin{center}
    \begin{tikzpicture}[scale=.8]
      \foreach \x in {0,5,10}{
      \fill(\x,0) circle (1);
      \fill[black!2] (\x,0) circle (.8);
      \fill (\x,0) circle (.1);
      \foreach\theta in {0,30,...,359}
      \draw[thick,rotate around={\theta:(\x,0)}] (\x,0)--+(1,0);
      }
      
      \foreach \y in {-1,0,1}{
        \fill[red] (0,\y) circle (.05);
        \draw[vectors,red] (0,\y)--(1.5,\y) node[right]{$\vec v_\text{cm}$};
      }
      \node[above] at (0,1.5) {Translation};
      
      \node at (3,0) {\Huge$+$};

      \begin{scope}[vectors,blue]
        \draw[rotate around={-40:(5,0)}] (5,.5) arc (90:0:.5)
        node[pos=1.2]{$\omega$};
        \draw (5,1)--(6.5,1) node[above]{$v=\omega R$};
        \draw (5,-1)--(3.5,-1) node[below]{$v=\omega R$};
      \end{scope}
      \foreach \y in {-1,1} \fill[blue] (5,\y) circle (.05);
      \node[above] at (5,1.5) {Rotation};
      \node at (7.5,0) {\Huge$=$};

      \begin{scope}[violet]
        \foreach \y in {0,1} \fill (10,\y) circle (.05);
        \draw[vectors] (10,1)--+(3,0) node[above]{$v=2\omega R$};
        \draw[vectors] (10,0)--+(1.5,0) node[right]{$v_\text{cm}=\omega R$};
        \fill (10,-1) circle (.05) node[below]{$v=0$};
      \end{scope}
    \end{tikzpicture}
  \end{center}
  At the contact point, if there is no slipping, the speed is zero. This can
  only occur when

  \eq{-.1in}{
    v_\text{cm}=\omega R\quad\longrightarrow\quad
    a_\text{cm}=\alpha R
  }
\end{frame}



\section{Pure Rolling Problem}

\begin{frame}{Pure Rolling Problems}
  \begin{columns}[T]
    \column{.35\textwidth}
    \begin{tikzpicture}[scale=1.1]
  \shade[ball color=gray!5] circle (1);
  \draw[thick] (-2,-1)--(2,-1);
  \draw[vectors] (1.1,0)--+(.7,0) node[right]{$\vec v_\text{cm}$};
  \draw[vectors] (0,1.2) arc (90:60:1.2) node[pos=1.15]{$\omega$};
  \draw[axes] (-2,-.75)--+(.5,0) node[right]{$x$};
  \draw[axes] (-2,-.75)--+(0,.5) node[above]{$y$};
  \fill circle (.05) node[right]{cm};
  \draw[vectors,blue] (0,0)--(0,-1.5) node[below]{$\vec F_g$};
  \draw[vectors,red] (-.03,-1)--(-.03,.5) node[above]{$\vec F_n$};
\end{tikzpicture}


    \column{.65\textwidth}
    In a \textbf{pure rolling} problem, a smooth solid sphere (any object that
    can roll with do) rolls along a smooth surface without slipping
    \begin{itemize}
    \item Assume both the sphere and the surface are both perfectly rigid
      (they do not deform)
    \item There are only two forces acting on the sphere:
      \begin{itemize}
      \item Gravitational force $\vec F_g$ (acting at the centre of mass)
      \item Normal force $\vec F_n$ (acting at the point of contact)
      \end{itemize}
    \item There is no friction
    \end{itemize}
  \end{columns}
\end{frame}



\begin{frame}{Pure Rolling Problem: Translational Motion}
  \begin{columns}
    \column{.65\textwidth}
    There is no net force, therefore the translational state
    %($\vec v_\text{cm}$)
    of the sphere is constant

    \eq{-.1in}{
      \sum\vec F=\vec 0\quad\longrightarrow\quad \vec v_\text{cm}=\text{constant}
    }

    The centre of mass of the sphere moves constant velocity and no acceleration
    \column{.35\textwidth}
    \begin{tikzpicture}[scale=1.1]
  \shade[ball color=gray!5] circle (1);
  \draw[thick] (-2,-1)--(2,-1);
  \draw[vectors] (1.1,0)--+(.7,0) node[right]{$\vec v_\text{cm}$};
  \draw[vectors] (0,1.2) arc (90:60:1.2) node[pos=1.15]{$\omega$};
  \draw[axes] (-2,-.75)--+(.5,0) node[right]{$x$};
  \draw[axes] (-2,-.75)--+(0,.5) node[above]{$y$};
  \fill circle (.05) node[right]{cm};
  \draw[vectors,blue] (0,0)--(0,-1.5) node[below]{$\vec F_g$};
  \draw[vectors,red] (-.03,-1)--(-.03,.5) node[above]{$\vec F_n$};
\end{tikzpicture}

  \end{columns}
\end{frame}




\begin{frame}{Pure Rolling Problem: Rotational Motion}
  \begin{columns}
    \column{.65\textwidth}
    The line of action of both normal force $\vec F_n$ and gravity $\vec F_g$
    go through the centre of mass, therefore neither force generate a torque
    about the CM. Since there is no net torque, the rotational state $\omega$
    is constant:
    
    \eq{-.1in}{
      \sum\tau=0\quad\longrightarrow \quad\omega=\text{constant}
    }

    In \emph{theory}, this sphere will roll along with angular speed $\omega$
    and speed $v_\text{cm}=\omega R$ forever

    \column{.35\textwidth}
    \begin{tikzpicture}[scale=1.1]
  \shade[ball color=gray!5] circle (1);
  \draw[thick] (-2,-1)--(2,-1);
  \draw[vectors] (1.1,0)--+(.7,0) node[right]{$\vec v_\text{cm}$};
  \draw[vectors] (0,1.2) arc (90:60:1.2) node[pos=1.15]{$\omega$};
  \draw[axes] (-2,-.75)--+(.5,0) node[right]{$x$};
  \draw[axes] (-2,-.75)--+(0,.5) node[above]{$y$};
  \fill circle (.05) node[right]{cm};
  \draw[vectors,blue] (0,0)--(0,-1.5) node[below]{$\vec F_g$};
  \draw[vectors,red] (-.03,-1)--(-.03,.5) node[above]{$\vec F_n$};
\end{tikzpicture}

  \end{columns}
\end{frame}



\begin{frame}{Reality: Rolling Resistance}
  In reality, the rolling sphere will slow down and eventually come to a stop,
  because \emph{nothing is perfectly rigid}: both the sphere and the surface
  deform when they make contact
  \begin{columns}
    \column{.7\textwidth}
    \begin{itemize}
    \item Example: a car's tires flatten when they make contact with the ground
    \item The normal force is larger in magnitude on the front side than on the
      other
    \item $\vec F_n$ exerts both a horizontal force to slow down the sphere, as
      well as a torque to slow down its rotation
    \item The normal force does not point towards the CM because of the
      deformation.
    \end{itemize}

    \column{.3\textwidth}
    \vspace{.3in}
    \pic1{OAGZy}
  \end{columns}
\end{frame}




\section{Rolling on Inclined Surface}


\begin{frame}{Rolling with Slipping on Rough Surface}
  \begin{columns}
    \column{.3\textwidth}
    \begin{tikzpicture}
  \shade[ball color=gray!5] circle (1);
  \draw[thick] (-2,-1)--(2,-1);
  \draw[vectors] (1.1,0)--+(.8,0) node[right]{$\vec v_\text{cm}$};
  \draw[vectors] (0,1.2) arc (90:60:1.2) node[pos=1.2]{$\omega$};

  \fill circle (.05) node[right]{cm};

  \draw[vectors,blue] (-.03,0)--(-.03,-1.5) node[below]{$\vec F_g$};
  \draw[vectors,red] (0,-1)--(0,.5) node[above]{$\vec F_n$};
  \draw[vectors,violet] (0,-1)--(1,-1) node[below]{$\vec f_k$};

  \draw[axes] (-2,-.75)--+(.5,0) node[right]{$x$};
  \draw[axes] (-2,-.75)--+(0,.5) node[above]{$y$};
\end{tikzpicture}


    \column{.7\textwidth}
    What if the rolling sphere is slipping against the surface? In this
    example, the sphere is rolling at rate $r\omega>v_\text{cm}$
    \begin{itemize}
    \item Slippage at the point of contact between the sphere and the flat
      surface
    \item There is \emph{kinetic} friction $f_k=\mu_kF_n$ in the $+x$
      direction (towards the right). The friction force generates:
      \begin{itemize}
      \item A net force $F_\text{net}$ in the $+\iii$ direction, towards the
        right
      \item A net torque $\tau_\text{net}$ in the $+\kkk$ direction, i.e.\
        counter clockwise
      \end{itemize}
    \end{itemize}
  \end{columns}
\end{frame}



\begin{frame}{Rolling with Slipping on Rough Surface: Translational Motion}
  \begin{columns}[T]
    \column{.3\textwidth}
    \begin{tikzpicture}
  \shade[ball color=gray!5] circle (1);
  \draw[thick] (-2,-1)--(2,-1);
  \draw[vectors] (1.1,0)--+(.8,0) node[right]{$\vec v_\text{cm}$};
  \draw[vectors] (0,1.2) arc (90:60:1.2) node[pos=1.2]{$\omega$};

  \fill circle (.05) node[right]{cm};

  \draw[vectors,blue] (-.03,0)--(-.03,-1.5) node[below]{$\vec F_g$};
  \draw[vectors,red] (0,-1)--(0,.5) node[above]{$\vec F_n$};
  \draw[vectors,violet] (0,-1)--(1,-1) node[below]{$\vec f_k$};

  \draw[axes] (-2,-.75)--+(.5,0) node[right]{$x$};
  \draw[axes] (-2,-.75)--+(0,.5) node[above]{$y$};
\end{tikzpicture}


    \column{.7\textwidth}
    The equation of motion for translational motion

    \vspace{-.35in}{\large
      \begin{align*}
        \sum F_x&=f_k=ma_\text{cm}\\
        \sum F_y&=F_n-mg=0
      \end{align*}
    }

    \vspace{-.02in}Kinetic friction $f_k$ causes the CM to accelerate to the
    right

    \eq{-.15in}{
      f_k=\mu_kmg=ma_\text{cm}
    }

    \vspace{-.15in}The acceleration of the centre of mass is constant:

    \eq{-.18in}{
      a_\text{cm}=\mu_kg
    }

    \vspace{-.15in}The velocity along the $x$ direction increases with time:

    \eq{-.18in}{
      v_\text{cm}(t)=v_0+\mu_kgt
    }
  \end{columns}
\end{frame}



\begin{frame}{Rolling with Slipping on Rough Surface}
  \begin{columns}[T]
    \column{.3\textwidth}
    \begin{tikzpicture}
  \shade[ball color=gray!5] circle (1);
  \draw[thick] (-2,-1)--(2,-1);
  \draw[vectors] (1.1,0)--+(.8,0) node[right]{$\vec v_\text{cm}$};
  \draw[vectors] (0,1.2) arc (90:60:1.2) node[pos=1.2]{$\omega$};

  \fill circle (.05) node[right]{cm};

  \draw[vectors,blue] (-.03,0)--(-.03,-1.5) node[below]{$\vec F_g$};
  \draw[vectors,red] (0,-1)--(0,.5) node[above]{$\vec F_n$};
  \draw[vectors,violet] (0,-1)--(1,-1) node[below]{$\vec f_k$};

  \draw[axes] (-2,-.75)--+(.5,0) node[right]{$x$};
  \draw[axes] (-2,-.75)--+(0,.5) node[above]{$y$};
\end{tikzpicture}


    \column{.7\textwidth}
    The constant kinetic friction force generates a constant net torque in the
    $+z$ (CCW) direction
    
    \eq{-.15in}{
      \tau_\text{net}=f_kR=I\alpha
    }

    \vspace{-.2in}Therefore there is a constant angular acceleration in the
    $+z$ direction:

    \eq{-.15in}{
      \alpha=\frac{\textcolor{orange}{f_k}R}{\color{violet}I}
      =\frac{\textcolor{orange}{\mu_kmg}R}{\color{violet}\frac25mR^2}
      =\frac{5\mu_kg}{2R}
    }

    \vspace{-.05in}Since angular acceleration (CCW) is in the opposite
    direction to angular velocity (CW), angular \emph{speed} $|\omega|$
    decreases over time:

    \eq{-.25in}{
      |\omega|=|\omega_0|-\alpha t
    }
  \end{columns}
\end{frame}



\begin{frame}{Rolling with Slipping on Rough Surface}
  \begin{columns}[T]
    \column{.3\textwidth}
    \begin{tikzpicture}
  \shade[ball color=gray!5] circle (1);
  \draw[thick] (-2,-1)--(2,-1);
  \draw[vectors] (1.1,0)--+(.8,0) node[right]{$\vec v_\text{cm}$};
  \draw[vectors] (0,1.2) arc (90:60:1.2) node[pos=1.2]{$\omega$};

  \fill circle (.05) node[right]{cm};

  \draw[vectors,blue] (-.03,0)--(-.03,-1.5) node[below]{$\vec F_g$};
  \draw[vectors,red] (0,-1)--(0,.5) node[above]{$\vec F_n$};
  \draw[vectors,violet] (0,-1)--(1,-1) node[below]{$\vec f_k$};

  \draw[axes] (-2,-.75)--+(.5,0) node[right]{$x$};
  \draw[axes] (-2,-.75)--+(0,.5) node[above]{$y$};
\end{tikzpicture}


    \column{.7\textwidth}
    \begin{itemize}
    \item Unlike the no-slip case where where angular acceleration is related
      to translational acceleration by the radius,
      %, i.e.\ $a=\alpha R$,
      in this case, there is \textbf{no relationship between $\alpha$ and $a$.}
    \item There is also no relationship between the velocity and the centre
      of mass and the angular velocity
    \item The speed of the sphere $v_\text{cm}$ increases, while the angular
      speed $\omega$ decreases, until\ldots
    \item When $v_\text{cm}=\omega R$, the sphere stops slipping, and the
      problem returns to the no-slip case
    \end{itemize}
  \end{columns}
\end{frame}



\section{Rolling on Inclined Surface}

\begin{frame}{Rolling on an Inclined Surface}
  \begin{columns}
    \column{.35\textwidth}
    \input{roll-down-ramp}

    \column{.65\textwidth}
    For a rigid and smooth sphere of radius $R$ rolling down a ramp of angle
    $\phi$ without slippage, three forces act on the sphere as it rolls down
    the ramp
    \begin{itemize}
    \item The weight ($\vec F_g$) of the sphere acts at the CM
    \item The normal force ($F_n$) acts at the point of contact
    \item The static friction ($f_s$) act at the point of contact
    \end{itemize}
    Only static friction generates a torque about the CM in the clockwise
    direction
    \begin{itemize}
    \item If $f_s$ is not present, there would have been nothing that causes
      the sphere to rotate.
    \end{itemize}
  \end{columns}
\end{frame}



\begin{frame}{Rolling on an Inclined Surface}
  \begin{columns}
    \column{.35\textwidth}
    \input{roll-down-ramp}

    \column{.65\textwidth}
    There are two equations of motion for the translational motion, along the
    $x$ and $y$ axes:
  
    \vspace{-.3in}{\large
      \begin{align*}
        \sum F_x&=mg\sin\phi-f_s=ma_\text{cm}\\
        \sum F_y&=F_n-mg\cos\phi=0
        %\sum\tau &=Rf_s=I_\text{cm}\alpha
      \end{align*}
    }

    And one equation of motion for the rotational motion. (The axis of
    rotation is along the $z$ axis.)

    \eq{-.1in}{
      \sum\tau =Rf_s=I_\text{cm}\alpha
    }
    
    %At this point, static friction $f_s$ is \emph{not} known. The coefficient
    %of static friction ($\mu_s$) only tells us the \emph{maximum} static
    %friction, not the \emph{actual} friction. (We will instead use it to check
    %if the answer makes sense.)
  \end{columns}
\end{frame}



\begin{frame}{Rolling on an Inclined Surface}
  \begin{columns}
    \column{.35\textwidth}
    \input{roll-down-ramp}

    \column{.65\textwidth}
    At this point, the magnitude of static friction $f_s$ is \emph{not} known.
    The coefficient of static friction ($\mu_s$) only tells us the
    \emph{maximum} static friction, not the \emph{actual} friction.

    \vspace{.1in}The fact that we have \emph{assumed} that the friction is
    static need to be checked after we have solved the problem. If the friction
    exceeds the  maximum allowable for static friction, the sphere will slip,
    and the problem becomes much more difficult.
  \end{columns}
\end{frame}





\begin{frame}{Rolling on an Inclined Surface}
  \begin{columns}
    \column{.35\textwidth}
    \input{roll-down-ramp}

    \column{.65\textwidth}
    For non-slip case, angular and translational acceleration are related using
    relative motion:

    \eq{-.15in}{
      a_\text{cm}=\alpha R
    }
    
    \vspace{-.22in}Solving for the static friction:

    \eq{-.1in}{
      f_s=\frac{I_\text{cm}\alpha}R=
      \frac25mR^2\cdot\frac aR\cdot\frac1R=\frac25ma_\text{cm}
    }

    \vspace{-.1in}It is substituted into the force equation in the $\iii$
    direction to solve for the acceleration of the CM down the ramp:

    \eq{-.1in}{
      mg\sin\phi-\frac25 ma_\text{cm}=ma_\text{cm}
    }
  \end{columns}
\end{frame}



\begin{frame}{Rolling on an Inclined Surface}
  \begin{columns}
    \column{.35\textwidth}
    \input{roll-down-ramp}

    \column{.65\textwidth}
    The acceleration of the centre of mass is constant in time:

    \eq{-.25in}{
      a_\text{cm}=\frac57 g\sin\phi
    }

    \vspace{-.1in}For comparison, an object \emph{sliding} without friction
    down the same ramp is higher than the pure-rolling case:
    
    \eq{-.22in}{
      a=g\sin\phi
    }
    
    \vspace{-.2in}If the sphere starts from rest, the speed of the sphere when
    it reaches the bottom of the ramp, a distance $d$ away, would be:

    \eq{-.2in}{
      v=\sqrt{2ad}=\sqrt{\frac{10}7gd\sin\phi}
    }
  \end{columns}
\end{frame}


\begin{frame}[t]{Sanity Check}
  Whenever we solve problems that have static friction, we must do a
  ``sanity check'' to make sure that the static friction has not exceeded the
  maximum value:

  \eq{-.1in}{
    f_s < \mu F_n
  }
  
  \vspace{1.3in}
  If this condition is met, then we will know for sure that the actual static
  friction is less than maximum, and the sphere will not slip.
\end{frame}


\section{Pulley Problems}

\begin{frame}{Pulley Problem}
  \begin{columns}
    \column{.25\textwidth}
    \centering
    \begin{tikzpicture}[scale=.7]
  \draw[ultra thick,brown] (-1,-1.6)--(-1,0);
  \draw[ultra thick,brown] (1,0)--(1,-3);
  \draw[thick,fill=gray] circle (1.05);
  \draw[thick,fill=gray!40] circle (.95);
  \draw[line width=5.5] (0,-.15)--(0,2);
  \draw[very thick] (-2,2)--(2,2);
  \fill[white] circle (.1);
  \draw[mass] (-1.5,-1.6) rectangle +(1,-1) node[midway]{$m_1$};
  \draw[mass] (.5,-3) rectangle +(1,-1.4) node[midway]{$m_2$};
\end{tikzpicture}   


    \column{.75\textwidth}
    Pulley problems like these are common in most dynamics problems (e.g.\ as
    we gave seen in Classes 2 and 3). However, when we solve these problems, we
    always insist that the pulleys are massless and frictionless.
    \begin{itemize}
    \item We can get away with this \emph{if} the mass of the pulley is small
      compared to $m_1$ and $m_2$
    \item We cannot treat the problem the same way if the mass of the pulley is
      large
    \end{itemize}
  \end{columns}
\end{frame}



\begin{frame}[t]{Pulley Problem}
  \begin{columns}
    \column{.25\textwidth}
    \centering
    \begin{tikzpicture}[scale=.7]
  \draw[ultra thick,brown] (-1,-1.6)--(-1,0);
  \draw[ultra thick,brown] (1,0)--(1,-3);
  \draw[thick,fill=gray] circle (1.05);
  \draw[thick,fill=gray!40] circle (.95);
  \draw[line width=5.5] (0,-.15)--(0,2);
  \draw[very thick] (-2,2)--(2,2);
  \fill[white] circle (.1);
  \draw[mass] (-1.5,-1.6) rectangle +(1,-1) node[midway]{$m_1$};
  \draw[mass] (.5,-3) rectangle +(1,-1.4) node[midway]{$m_2$};
\end{tikzpicture}   

    
    \column{.75\textwidth}
    Instead, we have to solve the problem using rotational dynamics:
    \begin{itemize}
    \item As the masses accelerate upwards/downwards, cable tension must also
      generate a torque on the pulley
    \item The tension will not be constant along the cable
    \item If the cable is wrapped tightly around the pulley, then we can also
      expect the relation to hold true

      \eq{-.1in}{
        v=\omega R\quad\quad a=\alpha R
      }
    \end{itemize}
  \end{columns}
\end{frame}



\begin{frame}[t]{Simple Example}
  \begin{columns}
    \column{.2\textwidth}
    \centering
    \begin{tikzpicture}[scale=.7]
      \draw[ultra thick,brown] (1,0)--(1,-3);
      \draw[thick,fill=gray] circle (1.05);
      \draw[thick,fill=gray!40] circle (.95);
      \draw[line width=5.5] (0,-.15)--(0,2);
      \draw[very thick] (-2,2)--(2,2);
      \fill[white] circle (.1) node[black,left=4]{$M$};
      \draw[mass] (.5,-3) rectangle +(1,-1.4) node[midway]{$m$};
    \end{tikzpicture}   

    \column{.8\textwidth}
    \textbf{Example:} A mass $m$ is connected to a cable that is wrapped
    tightly around pulley with radius $R$, mass $M$ and a rotational inertia of
    $I=\frac12MR^2$. The pulley can rotate about its axis without friction.
    \begin{enumerate}
    \item What is the magnitude of the acceleration of the mass?
    \item What is the magnitude of the tension force?
    \item If the mass is released from rest, what is its speed after falling
      a distance $d$?
    \end{enumerate}
  \end{columns}
  
\end{frame}



\begin{frame}[t]{A More Complicated Example}
  \begin{columns}[T]
    \column{.2\textwidth}
    \centering
    \begin{tikzpicture}[scale=.7]
  \draw[ultra thick,brown] (-1,-1.6)--(-1,0);
  \draw[ultra thick,brown] (1,0)--(1,-3);
  \draw[thick,fill=gray] circle (1.05);
  \draw[thick,fill=gray!40] circle (.95);
  \draw[line width=5.5] (0,-.15)--(0,2);
  \draw[very thick] (-2,2)--(2,2);
  \fill[white] circle (.1);
  \draw[mass] (-1.5,-1.6) rectangle +(1,-1) node[midway]{$m_1$};
  \draw[mass] (.5,-3) rectangle +(1,-1.4) node[midway]{$m_2$};
\end{tikzpicture}   


    \column{.8\textwidth}
    \textbf{Example:} Two masses, $m_1$ and $m_2$, are connected to each other,
    and wrapped around a pulley with radius $R$, mass $m_3$ and a rotational
    inertia of $I=\frac12m_3R^2$. The pulley can rotate about its axis without
    friction. Assuming that $m_1>m_2$,
    \begin{enumerate}
    \item What is the magnitude of the acceleration of the  masses?
    \item If the masses are released from rest, what is their speed after
      falling a distance $d$?
    \end{enumerate}
  \end{columns}
\end{frame}



\section{Energy Conservation in Rotational Motion}

\begin{frame}{Translational Work}

  Recall that translational work is done when a \emph{force} displaces an
  object by $\dl x$:

  \eq{-.1in}{
    W_\text{trans}=\int_{\mathcal C}\vec F(\vec x)\cdot\dl\vec x
  }

  Net translational work done on an object is equal to the change in
  translational kinetic energy:

  \eq{-.15in}{
    W_\text{net}=\Delta K_t
    \quad\text{\normalsize where}\quad
    K_t=\frac12mv_\text{cm}^2
  }

  %It does not matter what $\vec F_\text{net}$ consists of. If $\vec F_\text{net}$
  %does any work, it will \emph{always} be equal to the change in translational
  %kinetic energy.
\end{frame}



\begin{frame}{Rotational Work}

  Likewise, rotational work is done when a \emph{torque} displaces an object by
  $\dl\theta$ along the direction of rotation:

  \eq{-.2in}{
    W_\text{rot}=\int_{\theta_0}^{\theta_1}\tau\dl\theta
  }

  The net rotational work done on an object
  %(i.e.\ work done by the net torque)
  is equal to the change in rotational kinetic energy:

  \eq{-.15in}{
    W_\text{net,rot}=\Delta K_r
    \quad\text{\normalsize where}\quad
    K_r=\frac12I\omega^2
  }
  %
  %\vspace{-.1in}It does not matter what $\tau_\text{net}$ consists of; as long as
  %$\tau_\text{net}$ does work, it will \emph{always} be equal to the change in
  %rotational kinetic energy.
\end{frame}



\begin{frame}{Kinetic Energy of a Rotating System}
  The total kinetic energy of a rotating system is the sum of its translational
  and rotational kinetic energies at its centre of mass:

  \eq{-.1in}{
    \boxed{
      K_\text{tot}=K_t+K_r
      =\frac12mv_\text{cm}^2+\frac12I_\text{cm}\omega^2}
  }
  
  In this case, $I_\text{cm}$ is calculated at the CM. For simple
  problems, we only need to compute rotational kinetic energy at the pivot:
  
  \eq{-.1in}{
    \boxed{K_\text{rot}=\frac12I_\text{pivot}\omega^2}
  }
  
  In this case, the $I_\text{pivot}$ is calculated at the pivot.
  \textbf{IMPORTANT:} $I_\text{cm}\neq I_\text{pivot}$
\end{frame}



\begin{frame}{Conservative Forces}
  Forces that are related to a potential energy are called
  \textbf{conservative forces}. They include:
  \begin{itemize}
  \item Gravitational force $\vec F_g$
  \item Spring force $\vec F_s$
  \item Electrostatic force $\vec F_q$
  \item Magnetic force $\vec F_m$
  \item Nuclear forces
  \end{itemize}
\end{frame}




\begin{frame}{Conservative Forces}
  All conservative forces share these common properties:
  \begin{enumerate}
  \item The work done by these forces relate to a change of a potential energy.
    This is often called the \textbf{work-potential energy theorem}:

    \eq{-.1in}{
      \boxed{
        W_\text{conservative}=-\Delta U
      }
    }
    \begin{itemize}
    \item\vspace{-.17in}Positive work decreases this related potential energy
    \item Negative work increases this related potential energy
    \end{itemize}
  \item $+$ work done by a conservative force converts potential energy
    into kinetic energy
  \item $-$ work done by a conservative force converts kinetic energy
    into potential energy
  \item The work done by a conservative force is \emph{path independent}
    %(depends only on end points)
  \end{enumerate}
\end{frame}



\begin{frame}{Conservative Forces}
  By the fundamental theorem of calculus, any conservative forces $\vec F$
  must be the negative gradient of the potential energies:

  \eq{-.1in}{
    \boxed{
      \vec F=-\nabla U=
      -\diffp Ux\hat\imath-\diffp Uy\hat\jmath-\diffp Uz\hat k
    }
  }

  In one-dimension (which is likely to be what you may encounter in AP Physics
  C):

  \eq{-.1in}{
    \boxed{
      F=-\diff Ux
    }
  }

  The direction of a conservative force \emph{always} decreases the potential
  energy. (Pay attention to the negative sign.)
\end{frame}



\begin{frame}{Examples of Non-Conservative Force}
  The majority of forces are \textbf{non-conservative}. The common forces
  discussed in the previous class are generally non-conservative:
  \begin{itemize}
  \item Applied force
  \item Tension force
  \item Normal force
  \item Friction
  \item Drag (fluid resistance)
  \end{itemize}
  The work-energy theorem still applies for non-conservative forces
\end{frame}



\begin{frame}{Work by Non-Conservative Forces}
  The work done by non-conservative forces differs from conservative forces in
  that:
  \begin{itemize}
  \item There is \textbf{no related potential energies}: the work done by a
    non-conservative force transforms energy
    \begin{itemize}
    \item From kinetic energy of one object to another
    \item From one form of kinetic energy to another
    \item From kinetic energy into thermal energy
    \end{itemize}
  \item The work is generally \textbf{path dependent}
  \end{itemize}
\end{frame}




\section{Rolling on Inclined Surface}

\begin{frame}{Rolling on an Inclined Surface}
  We now re-examine the example problem for a rigid sphere of radius
  $R$ rolling down a ramp of angle  $\phi$ without slippage, this time using
  the conservation of energy.

  \vspace{.2in}
  \begin{columns}
    \column{.35\textwidth}
    \input{roll-down-ramp}

    \column{.65\textwidth}
    Three forces act on the sphere as it rolls down the ramp
    \begin{itemize}
    \item Gravitational force ($\vec F_g$) acts at the CM
    \item Normal force ($\vec F_n$) acts at the point of contact
    \item Static friction ($\vec f_s$) act at the point of contact
    \end{itemize}
    Only static friction generates a torque about the CM in the clockwise
    direction
  \end{columns}
\end{frame}



\begin{frame}{Rolling on an Inclined Surface}
  The CM translates down the ramp, along the $+x$ direction. The sphere will
  also rotate in the clockwise direction. The work done by the three forces
  acting on the sphere:
  \begin{columns}
    \column{.35\textwidth}
    \input{roll-down-ramp}

    \column{.65\textwidth}
    
    \begin{itemize}
    \item Gravitational force ($\vec F_g$)
      \begin{itemize}
      \item $+$ translational work
      \item No rotational work
      \end{itemize}
    \item Normal force ($\vec F_n$)
      \begin{itemize}
      \item No translational work
      \item No rotational work
      \end{itemize}
    \item Static friction ($\vec f_s$)
      \begin{itemize}
      \item $-$ translational work
      \item $+$ rotational work
      \end{itemize}
    \end{itemize}
  \end{columns}
\end{frame}



\begin{frame}{Rolling on an Inclined Surface}
  In this case, the system consists of the the sphere and Earth, like the
  (much simpler) system of a mass sliding down a ramp.

  \vspace{.2in}
  \begin{columns}
    \column{.35\textwidth}
    \input{roll-down-ramp}

    \column{.65\textwidth}
    The total energy of the system therefore includes:
    \begin{itemize}
    \item The translational kinetic energy of the sphere
    \item The rotational kinetic energy of the sphere
    \item The gravitational potential energy stored between the sphere \& Earth
    \end{itemize}
    The system is an isolated system, because there is no external work done
    \begin{itemize}
    \item Although friction is non-conservative, it is \emph{internal} to the
      system
    \end{itemize}
  \end{columns}
\end{frame}



\begin{frame}{Rolling on an Inclined Surface}
  \begin{columns}
    \column{.35\textwidth}
    \input{roll-down-ramp}

    \column{.65\textwidth}
    As the sphere rolls downward without slipping, the total mechanical
    energy is conserved, i.e.:

    \eq{-.1in}{
      \boxed{
        K_t + K_r + U_g = \text{constant}
      }
    }
    \begin{itemize}
    \item\vspace{-.1in}The positive translational work done by gravity converts
      $U_g$ into $K_t$
    \item The work done by static friction converts $K_t$ to $K_r$
    \item All the energy stays inside the system, therefore the system is an
      isolated system
    \end{itemize}
  \end{columns}
\end{frame} 



\begin{frame}[t]{Back to the Simple Pulley Example}
  \begin{columns}
    \column{.2\textwidth}
    \centering
    \begin{tikzpicture}[scale=.7]
      \draw[ultra thick,brown] (1,0)--(1,-3);
      \draw[thick,fill=gray] circle (1.05);
      \draw[thick,fill=gray!40] circle (.95);
      \draw[line width=5.5] (0,-.15)--(0,2);
      \draw[very thick] (-2,2)--(2,2);
      \fill[white] circle (.1) node[black,left=4]{$M$};
      \draw[mass] (.5,-3) rectangle +(1,-1.4) node[midway]{$m$};
    \end{tikzpicture}   

    \column{.8\textwidth}
    \textbf{Example:} A mass $m$ is connected to a cable that is wrapped
    tightly around pulley with radius $R$, mass $M$ and a rotational inertia of
    $I=\frac12MR^2$. The pulley can rotate about its axis without friction.
    If the mass is released from rest, what is its speed after falling
    a distance $d$?
  \end{columns}
\end{frame}



\begin{frame}{Back to the More Complicated Pulley Example}
  \begin{columns}[T]
    \column{.2\textwidth}
    \centering
    \begin{tikzpicture}[scale=.7]
  \draw[ultra thick,brown] (-1,-1.6)--(-1,0);
  \draw[ultra thick,brown] (1,0)--(1,-3);
  \draw[thick,fill=gray] circle (1.05);
  \draw[thick,fill=gray!40] circle (.95);
  \draw[line width=5.5] (0,-.15)--(0,2);
  \draw[very thick] (-2,2)--(2,2);
  \fill[white] circle (.1);
  \draw[mass] (-1.5,-1.6) rectangle +(1,-1) node[midway]{$m_1$};
  \draw[mass] (.5,-3) rectangle +(1,-1.4) node[midway]{$m_2$};
\end{tikzpicture}   


    \column{.8\textwidth}
    \textbf{Example:} Two masses, $m_1$ and $m_2$, are connected to each other,
    and wrapped around a pulley with radius $R$, mass $m_3$ and a rotational
    inertia of $I=\frac12m_3R^2$. The pulley can rotate about its axis without
    friction. Assuming that $m_1>m_2$, if the masses are released from rest,
    what is their speed after falling a distance $d$?
  \end{columns}
\end{frame}

\end{document}
